\documentclass[12pt]{article}

% === Кодировки и язык ===
\usepackage{cmap}
\usepackage[T2A]{fontenc}
\usepackage{mathtext}
\usepackage[utf8]{inputenc}
\usepackage[english, russian]{babel}

% === Цвета ===
\usepackage[dvipsnames]{xcolor}

% === Математика ===
\usepackage{amsmath, amssymb, wasysym, amsthm}
\usepackage{stmaryrd}
\usepackage{proof}

% === Графика и float'ы ===
\usepackage{wrapfig, float}
\usepackage{graphicx}

% === TikZ / PGFPlots ===
\usepackage{pgfplots}
\pgfplotsset{compat=newest}
\usepackage{pgfplotstable}
\usepackage{tikz, tikz-3dplot}
\usetikzlibrary{calc, patterns, angles, quotes, intersections, automata, positioning, arrows.meta}

% === Таблицы ===
\usepackage{tabularx}
\usepackage{xltabular}
\renewcommand\tabularxcolumn[1]{m{#1}}

% === Разметка ===
\usepackage{geometry}
\geometry{margin=0.5in}

% === Списки и прочее ===
\usepackage[shortlabels]{enumitem}
\usepackage{verbatim}

% ==========================================
% === ТЕОРЕМНЫЕ ОКРУЖЕНИЯ (amsthm) ===
% ==========================================
\theoremstyle{plain}
\newtheorem{theorem}{Теорема}[section]
\newtheorem{corollary}{Следствие}[section]
\newtheorem{lemma}{Лемма}[section]
\newtheorem{statement}{Утверждение}[section]
\newtheorem{axiom}{Аксиома}[section]
\newtheorem{law}{Закон}[section]

\theoremstyle{definition}
\newtheorem{definition}{Определение}[section]
\newtheorem{example}{Пример}[section]
\newtheorem{problem}{Задача}
\newtheorem{method}{Метод}[section]
\newtheorem{event}{Событие}[section]

\theoremstyle{remark}
\newtheorem{note}{Заметка}[section]
\newtheorem{property}{Свойство}[section]
\newtheorem{properties}{Свойства}[section]
\newtheorem{principle}{Правило}[section]
\newtheorem{person}{Личность}[section]
\newtheorem*{solution}{Решение}

% ==========================================
% === КРАСИВЫЕ РАМКИ (tcolorbox) ===
% ==========================================
\usepackage{tcolorbox}
\tcbuselibrary{skins, breakable, theorems}

% Базовые стили для разных типов блоков
\tcbset{
    % Стиль для теорем, лемм (Синий - строгий)
    thmstyle/.style={
        enhanced, breakable,
        frame hidden, boxrule=0pt,
        colback=blue!5!white, % Фон: 5% синего, 95% белого
        borderline west={2.5pt}{0pt}{blue!70!black}, % Вертикальная линия слева
        before skip=10pt, after skip=10pt,
        left=4pt, right=4pt, top=4pt, bottom=4pt
    },
    % Стиль для доказательств (фиолетовый - процесс рассуждения)
    proofstyle/.style={
        enhanced, breakable,
        frame hidden, boxrule=0pt,
        colback=Plum!5!white,
        borderline west={2.5pt}{0pt}{Plum!80!black},
        before skip=10pt, after skip=10pt,
        left=4pt, right=4pt, top=4pt, bottom=4pt,
        sharp corners
    },
    % Стиль для определений и аксиом (Зеленый - фундаментальный)
    defstyle/.style={
        enhanced, breakable,
        frame hidden, boxrule=0pt,
        colback=ForestGreen!5!white,
        borderline west={2.5pt}{0pt}{ForestGreen!80!black},
        before skip=10pt, after skip=10pt,
        left=4pt, right=4pt, top=4pt, bottom=4pt
    },
    % Стиль для примеров и задач (Оранжевый - практический)
    exstyle/.style={
        enhanced, breakable,
        frame hidden, boxrule=0pt,
        colback=orange!5!white,
        borderline west={2.5pt}{0pt}{orange!90!black},
        before skip=10pt, after skip=10pt,
        left=4pt, right=4pt, top=4pt, bottom=4pt
    },
    % Стиль для заметок и свойств (Серый - нейтральный)
    remstyle/.style={
        enhanced, breakable,
        frame hidden, boxrule=0pt,
        colback=gray!5!white,
        borderline west={2.5pt}{0pt}{gray!70!black},
        before skip=10pt, after skip=10pt,
        left=4pt, right=4pt, top=4pt, bottom=4pt
    }
}

% Применяем стили к существующим окружениям amsthm
% 1. Строгая математика
\tcolorboxenvironment{theorem}{thmstyle}
\tcolorboxenvironment{lemma}{thmstyle}
\tcolorboxenvironment{corollary}{thmstyle}
\tcolorboxenvironment{statement}{thmstyle}
\tcolorboxenvironment{law}{thmstyle}

% 2. Доказательство
\tcolorboxenvironment{proof}{proofstyle}
\renewcommand{\qed}{}

% 3. Определения
\tcolorboxenvironment{definition}{defstyle}
\tcolorboxenvironment{axiom}{defstyle}

% 4. Практика
\tcolorboxenvironment{example}{exstyle}
\tcolorboxenvironment{problem}{exstyle}
\tcolorboxenvironment{method}{exstyle}
\tcolorboxenvironment{event}{exstyle}

% 5. Заметки и прочее
\tcolorboxenvironment{note}{remstyle}
\tcolorboxenvironment{property}{remstyle}
\tcolorboxenvironment{properties}{remstyle}
\tcolorboxenvironment{principle}{remstyle}
\tcolorboxenvironment{person}{remstyle}
\tcolorboxenvironment{solution}{remstyle}

% === Гиперссылки (почти последним!) ===
\usepackage[
colorlinks=true,
linkcolor=blue,
urlcolor=blue,
citecolor=blue,
bookmarks=true,
bookmarksopen=false
]{hyperref}


\tikzset{
	dot/.style={circle, fill=black, inner sep=1pt},
	invisible/.style={dashed, thin, opacity=0.5},
	visible/.style={thick},
	main/.style={blue!70!black, thick},
	aux/.style={red!70!black, thin},
	plane/.style={fill=green!20, opacity=0.4},
	secant/.style={green!50!black, thick}
}

\begin{document}
	\tableofcontents
	\setcounter{tocdepth}{3}
	\newpage
	
	\section{Тетраэдр}
	
	\subsection{Медианы, высоты и средние линии}
	
	\begin{theorem}
		Если в тетраэдре две пары перпендикулярных скрещивающихся ребер, то он ортоцентрический (высоты пересекаются в одной точке).
		
		\begin{center}
			\begin{tikzpicture}[scale=1.2, line join=round, line cap=round]
				% Координаты вершин
				\coordinate (D) at (2, 4);
				\coordinate (A) at (0, 0);
				\coordinate (B) at (5, 0.5);
				\coordinate (C) at (1.5, 1.5);
				
				% Основание и скрытые линии
				\draw[invisible] (A) -- (C);
				\draw[invisible] (B) -- (C);
				\draw[invisible] (D) -- (C);
				
				% Видимые ребра
				\draw[visible] (A) -- (B) -- (D) -- cycle;
				
				% Точка пересечения высот (ортоцентр) H
				\coordinate (H) at (1.9, 1.4);
				\node[dot, label=right:$H$] at (H) {};
				
				% Высота из D на плоскость ABC
				\coordinate (Hd) at (1.9, 0.6);
				\draw[aux, dashed] (D) -- (H) -- (Hd);
				\node[below] at (Hd) {\small $H_D$};
				
				% Высота из A на плоскость BDC
				\coordinate (Ha) at (2.8, 2);
				\draw[aux, dashed] (A) -- (H) -- (Ha);
				\node[right] at (Ha) {\small $H_A$};
				
				% Подписи вершин
				\node[below left] at (A) {$A$};
				\node[below right] at (B) {$B$};
				\node[above] at (C) {$C$};
				\node[above] at (D) {$D$};
			\end{tikzpicture}
		\end{center}
	\end{theorem}
	
	\begin{theorem}
		Медианы тетраэдра пересекаются в одной точке и делятся в отношение $3 : 1$.
		
		\begin{center}
			\begin{tikzpicture}[scale=1.2, line join=round, line cap=round]
				\coordinate (D) at (2, 4);
				\coordinate (A) at (0, 0);
				\coordinate (B) at (5, 0);
				\coordinate (C) at (1.5, 1.5);
				
				% Скрытые ребра
				\draw[invisible] (A) -- (C);
				\draw[invisible] (B) -- (C);
				\draw[invisible] (D) -- (C);
				
				% Видимые ребра
				\draw[visible] (A) -- (B) -- (D) -- cycle;
				
				% Центроид грани ABC (G_D)
				\coordinate (Gd) at (2.16, 0.5); 
				% Центроид грани BCD (G_A)
				\coordinate (Ga) at (2.83, 1.83);
				
				% Точка пересечения медиан (Центроид тетраэдра M)
				% M = (A+B+C+D)/4. Gd = (A+B+C)/3. M = 3/4 Gd + 1/4 D.
				\coordinate (M) at (2.12, 1.37);
				
				% Рисуем медианы
				\draw[aux, dashed] (D) -- (M) -- (Gd);
				\draw[aux, dashed] (A) -- (M) -- (Ga);
				
				\node[dot, label=right:$M$] at (M) {};
				\node[below] at (Gd) {\small $G_D$};
				\node[right] at (Ga) {\small $G_A$};
				
				\node[below left] at (A) {$A$};
				\node[below right] at (B) {$B$};
				\node[above] at (C) {$C$};
				\node[above] at (D) {$D$};
				
				% Обозначение отношения
				\node[left, red] at (2.1, 2.7) {$3x$};
				\node[left, red] at (2.15, 0.9) {$x$};
			\end{tikzpicture}
		\end{center}
	\end{theorem}
	
	\begin{theorem}
		Средние линии тетраэдра пересекаются в одной точке и делятся в точке пересечения пополам.
		
		\begin{center}
			\begin{tikzpicture}[scale=1.2, line join=round, line cap=round]
				\coordinate (D) at (2, 4);
				\coordinate (A) at (0, 0);
				\coordinate (B) at (5, 0.5);
				\coordinate (C) at (1.5, 1.5);
				
				\draw[invisible] (A) -- (C);
				\draw[invisible] (B) -- (C);
				\draw[invisible] (D) -- (C);
				\draw[visible] (A) -- (B) -- (D) -- cycle;
				
				% Середины ребер
				\coordinate (M_AB) at (2.5, 0.25);
				\coordinate (M_CD) at (1.75, 2.75);
				\coordinate (M_AD) at (1, 2);
				\coordinate (M_BC) at (3.25, 1);
				
				% Точка пересечения
				\coordinate (M) at (2.125, 1.5);
				\node[dot, label=right:$M$] at (M) {};
				
				% Средние линии
				\draw[aux, dashed] (M_AB) -- (M_CD);
				\draw[aux, dashed] (M_AD) -- (M_BC);
				
				% Узлы для середин
				\node[dot, scale=0.5] at (M_AB) {};
				\node[dot, scale=0.5] at (M_CD) {};
				\node[dot, scale=0.5] at (M_AD) {};
				\node[dot, scale=0.5] at (M_BC) {};
				
				\node[below left] at (A) {$A$};
				\node[below right] at (B) {$B$};
				\node[above] at (C) {$C$};
				\node[above] at (D) {$D$};
			\end{tikzpicture}
		\end{center}
	\end{theorem}
	
	\subsection{Теорема Менелая}
	
	\begin{theorem}[Теорема Менелая для тетраэдра]
		Возьмем три точки $M, N, K$, лежащие в одной плоскости; $M \in AD$, $N \in BD$, $K \in AC$. Тогда если эта плоскость пересекает $BC$ в точке $P$, то:
		$$\dfrac{DM}{MA} \cdot \dfrac{AK}{KC} \cdot \dfrac{CP}{PB} \cdot \dfrac{BN}{ND} = 1$$
		
		\begin{center}
			\begin{tikzpicture}[scale=1.5, line join=round, line cap=round]
				% 1. Координаты вершин тетраэдра
				\coordinate (D) at (2, 4);
				\coordinate (A) at (0, 0);
				\coordinate (B) at (5, 0.5);
				\coordinate (C) at (1.5, 1.5);
				
				% 2. Точки на ребрах через интерполяцию (calc)
				% Точки M и N на передних ребрах
				\coordinate (M) at ($(A)!0.4!(D)$); 
				\coordinate (N) at ($(B)!0.6!(D)$);
				% Точки K и P на задних ребрах
				\coordinate (K) at ($(A)!0.65!(C)$);
				\coordinate (P) at ($(B)!0.4!(C)$); 
				
				% 3. Отрисовка невидимых ребер тетраэдра
				\draw[invisible] (A) -- (C);
				\draw[invisible] (B) -- (C);
				\draw[invisible] (D) -- (C);
				
				% 4. Заливка плоскости сечения
				\fill[plane] (M) -- (N) -- (P) -- (K) -- cycle;
				
				% 5. Отрисовка линий сечения
				% Линия MN видима (на передней грани ADB)
				\draw[secant] (M) -- (N);
				% Остальные линии сечения уходят на невидимые грани
				\draw[secant, dashed] (N) -- (P);
				\draw[secant, dashed] (P) -- (K);
				\draw[secant, dashed] (K) -- (M);
				
				% 6. Отрисовка видимых ребер тетраэдра
				\draw[visible] (A) -- (B) -- (D) -- cycle;
				
				% 7. Подписи и точки
				\node[below left] at (A) {$A$};
				\node[below right] at (B) {$B$};
				\node[above] at (C) {$C$};
				\node[above] at (D) {$D$};
				
				\node[dot, label=left:$M$] at (M) {};
				\node[dot, label=right:$N$] at (N) {};
				\node[dot, label=above left:$K$] at (K) {};
				\node[dot, label=below:$P$] at (P) {};
			\end{tikzpicture}
		\end{center}
	\end{theorem}
	
	\begin{statement}
		В ортоцентрическом тетраэдре суммы квадратов противоположных ребер равны:
		$$AB^2 + CD^2 = AC^2 + BD^2 = AD^2 + BC^2$$
		
		\begin{center}
			\begin{tikzpicture}[scale=1.5, line join=round, line cap=round]
				% Стили для параллелепипеда
				\tikzset{box/.style={gray, thin, opacity=0.4}}
				
				% Координаты вершин параллелепипеда (основание)
				\coordinate (V0) at (0, 0);       % Передний-левый-низ
				\coordinate (V1) at (3.5, 0.5);   % Передний-правый-низ
				\coordinate (V2) at (1.2, 1.2);   % Задний-левый-низ
				\coordinate (V3) at (4.7, 1.7);   % Задний-правый-низ (точка B)
				
				% Координаты вершин параллелепипеда (верх)
				\coordinate (V4) at (0, 3.5);     % Передний-левый-верх
				\coordinate (V5) at (3.5, 4.0);   % Передний-правый-верх (точка D)
				\coordinate (V6) at (1.2, 4.7);   % Задний-левый-верх (точка C)
				\coordinate (V7) at (4.7, 5.2);   % Задний-правый-верх
				
				% Вершины тетраэдра
				\coordinate (A) at (V0);
				\coordinate (B) at (V3);
				\coordinate (C) at (V6);
				\coordinate (D) at (V5);
				
				% 1. Отрисовка невидимых линий параллелепипеда
				\draw[box, dashed] (V2) -- (V0);
				\draw[box, dashed] (V2) -- (V3);
				\draw[box, dashed] (V2) -- (V6);
				
				
				% 2. Отрисовка видимых линий параллелепипеда
				\draw[box] (V0) -- (V1) -- (V3);
				\draw[box] (V4) -- (V5) -- (V7) -- (V6) -- cycle;
				\draw[box] (V0) -- (V4);
				\draw[box] (V1) -- (V5);
				\draw[box] (V3) -- (V7);
				
				% 3. Отрисовка ребер тетраэдра
				\draw[visible] (A) -- (B);
				\draw[visible] (A) -- (D);
				\draw[visible] (B) -- (D);
				\draw[visible] (C) -- (D);
				\draw[visible] (B) -- (C);
				\draw[visible] (A) -- (C);
				
				% 4. Подписи вершин тетраэдра
				\node[dot, label=below left:$A$] at (A) {};
				\node[dot, label=right:$B$] at (B) {};
				\node[dot, label=above:$C$] at (C) {};
				\node[dot, label=above:$D$] at (D) {};
				
				% Обозначения ребер параллелепипеда (для доказательства)
				\node[below, gray] at ($(V0)!0.5!(V1)$) {\small $x$};
				\node[right, gray] at ($(V1)!0.5!(V3)$) {\small $y$};
				\node[left, gray] at ($(V0)!0.5!(V4)$) {\small $z$};
			\end{tikzpicture}
		\end{center}
	\end{statement}

    \section{Координаты}

    \begin{definition}[Коллинеарные векторы]
        Лежат на параллельных прямых или на одной прямой.
    \end{definition}

    \begin{definition}[Равные векторы]
        Векторы, равные по модулю и направлению.
    \end{definition}

    \begin{definition}[Компланарные векторы]
        Векторы компланарные, если равные им векторы, имеющие общее начало, лежат в одной плоскости.
    \end{definition}

    \begin{note}
        Ось z~--- ось аппликат.
    \end{note}

    \begin{definition}[Расстояние]
        Расстояние между точками $A = (x_1, y_1, z_1)$ и $B = (x_2, y_2, z_2)$ равно $AB = \sqrt{(x_1 - x_2)^2 + (y_1 - y_2)^2 + (z_1 - z_2)^2}$.
    \end{definition}

    \begin{statement}[Деление отрезка в соотношении $\dfrac{m}{n}$]
        $\left( \dfrac{m x_1 + n x_2}{m + n}, \dfrac{m y_1 + n y_2}{m + n}, \dfrac{m z_1 + n z_2}{m + n} \right)$
    \end{statement}
    
    \subsection{Уравнение плоскости}
	
	\begin{statement}[Уравнение плоскости]
        Рассмотрим плоскость $\alpha$ в трехмерном пространстве. Зададим ее следующими параметрами:
        \begin{itemize}
            \item Вектор нормали $\vec{n}(a, b, c)$, такой что $\vec{n} \perp \alpha$.
            \item Фиксированная точка $M_0(x_0, y_0, z_0)$, принадлежащая плоскости ($M_0 \in \alpha$).
            \item Произвольная (текущая) точка плоскости $M(x, y, z) \in \alpha$.
        \end{itemize}
        Тогда уравнение плоскости такое:
        \[
        ax + by + cz + d = 0
        \]
    \end{statement}
    
    \begin{proof}[Вывод]
    	Так как точки $M_0$ и $M$ лежат в плоскости $\alpha$, вектор $\vec{M_0M}$ также лежит в этой плоскости. Поскольку $\vec{n} \perp \alpha$, то вектор нормали перпендикулярен любому вектору в плоскости:
        \[
        \vec{M_0M} \perp \vec{n} \implies \vec{M_0M} \cdot \vec{n} = 0
        \]
        
        Координаты вектора $\vec{M_0M} = (x - x_0, y - y_0, z - z_0)$. Записывая скалярное произведение в координатной форме, получаем \textbf{уравнение плоскости, проходящей через точку перпендикулярно вектору}:
        \[
            \boxed{a(x - x_0) + b(y - y_0) + c(z - z_0) = 0}
        \]
        
        Раскрыв скобки и обозначив свободный член как $d = -(ax_0 + by_0 + cz_0)$, мы переходим к \textbf{общему уравнению плоскости} (в конспекте названо каноническим видом):
        \[
            ax + by + cz + d = 0
        \]
    \end{proof}

    \subsubsection{Пример решения задачи}
    
    \begin{problem}
    	Найти уравнение плоскости, проходящей через три точки: $M_1(1, 2, -1)$, $M_2(-1, -1, 3)$ и $M_3(2, 3, -2)$.
    \end{problem}
    
    \begin{solution}
    	\begin{enumerate}
			\item Составим векторы:
       		$\vec{M_1M_2} = (-2, -3, 4)$
       		$\vec{M_1M_3} = (1, 1, -1)$
			\item Найдем вектор нормали $\vec{n}$ как векторное произведение $\vec{M_1M_2} \times \vec{M_1M_3}$ (или через систему уравнений, как в черновике). В результате расчетов получен вектор $\vec{n}(1, -2, -1)$.
			\item Подставим координаты вектора $\vec{n}$ и точки $M_1$ в уравнение (1):
       		\[
				1(x - 1) - 2(y - 2) - 1(z + 1) = 0
			\]
       		\[
				x - 1 - 2y + 4 - z - 1 = 0 \implies x - 2y - z + 2 = 0
			\]
		\end{enumerate}
    \end{solution}
    
    \subsection{Расстояние от точки до плоскости}
    
    \begin{problem}
    	Найти расстояние $\rho(M_0, \alpha)$ от точки $M_0(x_0, y_0, z_0)$ до плоскости $\alpha: ax + by + cz + d = 0$.
    \end{problem}
    
    \begin{solution}
        Пусть $M_1(x_1, y_1, z_1)$ — проекция точки $M_0$ на плоскость $\alpha$. Тогда вектор $\vec{M_1M_0}$ коллинеарен вектору нормали $\vec{n}(a, b, c)$:
        \[
        \vec{M_1M_0} = k \cdot \vec{n} \Rightarrow 
        \begin{cases} 
            x_0 - x_1 = ka \\
            y_0 - y_1 = kb \\
            z_0 - z_1 = kc
        \end{cases}
        \]
        Отсюда координаты точки $M_1$ выражаются как: $x_1 = x_0 - ka$, $y_1 = y_0 - kb$, $z_1 = z_0 - kc$.
        Так как $M_1 \in \alpha$, подставим эти координаты в уравнение плоскости:
        \[
        a(x_0 - ka) + b(y_0 - kb) + c(z_0 - kc) + d = 0
        \]
        \[
        ax_0 + by_0 + cz_0 + d = k(a^2 + b^2 + c^2) \Rightarrow k = \frac{ax_0 + by_0 + cz_0 + d}{a^2 + b^2 + c^2}
        \]
        
        Расстояние $\rho = |\vec{M_1M_0}| = \sqrt{(ka)^2 + (kb)^2 + (kc)^2} = |k|\sqrt{a^2 + b^2 + c^2}$.
        Подставляя значение $k$, получаем итоговую формулу:
        \[
            \boxed{\rho = \frac{|ax_0 + by_0 + cz_0 + d|}{\sqrt{a^2 + b^2 + c^2}}}
        \]
    \end{solution}

	\subsection{Уравнение прямой}
	\begin{statement}
		Прямая в пространстве может быть задана как линия пересечения двух непараллельных плоскостей:
    	\[
    		\begin{cases}
        		a_1x + b_1y + c_1z + d_1 = 0 \\
        		a_2x + b_2y + c_2z + d_2 = 0
    		\end{cases}
    	\]
		Но это не удобно. Удобнее задавать как
		\[
			\frac{x - x_1}{x_2 - x_1} = \frac{y - y_1}{y_2 - y_1} = \frac{z - z_1}{z_2 - z_1} = t
		\]
		Также есть параметрический вид:
		\[
			\begin{cases}
				x = x_1 + (x_2 - x_1)t \\
				y = y_1 + (y_2 - y_1)t \\
				z = z_1 + (z_2 - z_1)t
			\end{cases}
		\]
	\end{statement}
\end{document}
